\subsection{A clock reconstructed from the scalar records}
\label{sec:source-scalar-clock}

The finite execution also determines an action-time readout without using
the recorded step labels to assign its configuration layers. This construction
retains the supplied positive mass inner product, operator $\mathsf A$,
resource-to-site map and action normalization. It concerns one duration
shared by an entire configuration layer, not independent lapses at each site.

\paragraph{Recovering the layers.}
Authenticate each immutable write and every consumed resource, version,
writer and value. Draw directed edges from consumed writers to their readers.
The serial audit parent remains authenticated but is not a field input in
this action. Define roots to have height zero and every other vertex to have
height one plus its greatest parent height. An update reads its own field
two layers earlier and the nonempty current stencil, including itself.
Induction therefore gives height $j$ to every update of $u_j$, $1\le j\le21$.
No event identifier component or grouping of serialized positions enters
this definition.

Both seed configurations have height zero. Their roles are distinguished
in this particular graph by their numbers of height-one children: every
older-memory root has one, whereas each current-field root has four,
five, six or seven. The respective counts are $64,8,24,24,8$. With the
declared resource-to-site map, each recovered layer contains exactly one
write per site; the verifier checks alternating buffers, consecutive
versions and the full adjacent-layer action stencil. Thus all $23$ layers
are recovered in each trace. The two root sets remain causally incomparable;
the older seed is integrator memory, not a negative-time physical interval.
The complete retained read graph is required. For two symmetric roots feeding
a single update, graph symmetry prevents this seed-role distinction.

\begin{proposition}[Unique positive action durations]
\label{prop:source-scalar-clock-unique}
Let $\mathsf A$ be positive and self-adjoint in the supplied $M$ inner product.
For fixed endpoints consider the configuration-stationarity equations of
\begin{equation}
 \mathcal S_h[u]=\sum_{j=0}^{20}\left[
  \frac{\|u_{j+1}-u_j\|_M^2}{2h_j}
  -\frac{h_j}{4}\bigl(E(u_j)+E(u_{j+1})\bigr)\right],
 \qquad E(u)=\langle u,\mathsf A u\rangle_M,
 \label{eq:source-scalar-clock-action}
\end{equation}
where the $h_j>0$ are fixed but unknown parameters. At an interior center put
$x=u_j-u_{j-1}$, $y=\mathsf A u_j$, $z=u_{j+1}-u_j$ and suppose
\[
 D=\langle x,x\rangle_M\langle y,y\rangle_M
       -\langle x,y\rangle_M^2>0.
\]
If the full vector equation $z=\alpha x+\beta y$ holds, positive adjacent
durations exist exactly when $\alpha>0$ and $\beta<0$. They are unique:
\begin{equation}
 h_{j-1}^2=\frac{-2\beta}{\alpha(\alpha+1)},\qquad
 h_j=\alpha h_{j-1}.
 \label{eq:source-scalar-clock-duration}
\end{equation}
If every actual center $j=1,\ldots,20$ has $\alpha=1$, $\beta=-s$ with
$s>0$, all $21$ positive durations equal $\sqrt{s}$.
\end{proposition}
\begin{proof}
The discrete variational framework is supplied~\cite{marsdenWest}.
Positive Gram determinant implies independence of $x,y$. Configuration
variation in \eqref{eq:source-scalar-clock-action} gives
\[
 z/h_j-x/h_{j-1}+(h_j+h_{j-1})y/2=0.
\]
Its independent coefficients imply $h_j/h_{j-1}=\alpha$ and
$\beta=-h_j(h_j+h_{j-1})/2$, proving necessity and uniqueness.
Substitution proves sufficiency. Shared intervals must agree between
adjacent centers; for the stated uniform coefficients they do, including
the endpoint intervals determined by centers $1$ and $20$.
\end{proof}

The coefficients are obtained by the inverse Gram matrix, and their full
vector residual is separately checked; projection alone would miss a
component orthogonal to $x,y$. Exact replay of both intervention schedules
gives $40$ positive determinants, all exceeding $2083218601/10^{12}$,
and $(\alpha,\beta)=(1,-1/245)$. Hence the positive tick is $\sqrt5/35$
and the elapsed time from an assigned origin is $3/\sqrt5$. The zero
baseline has no independent pair and defines no such clock. Six Lean lemmas
formalize coefficient extraction, duration uniqueness and the Gram bridge;
an additional theorem composes these implications. The variation and the
actual finite-record values are established analytically and by exact replay.

At equal $h_j=\tau$, \eqref{eq:source-scalar-clock-action} differs from
the left-endpoint action yielding \eqref{eq:source-scalar-executed-update}
by the endpoint term
$\tau(E(u_0)-E(u_{21}))/4$, leaving fixed-endpoint stationarity unchanged.
Its variable-duration extension is a supplied quadrature choice. Durations
are not varied in deriving the configuration equation: unconstrained
variation in $h_j$ would give the additional, generally nonzero derivative
$-\|u_{j+1}-u_j\|_M^2/(2h_j^2)-(E(u_j)+E(u_{j+1}))/4$.

\begin{proposition}[Stable duration identification]
\label{prop:source-scalar-clock-stability}
In the same real Hilbert space let $F_0=[x\ y]$, $z=F_0c_0$,
$c_0=(\alpha_0,\beta_0)$, and $D>0$. Choose positive bounds
\[
 \ell_\alpha\ge\sqrt{\langle y,y\rangle_M/D},\qquad
 \ell_\beta\ge\sqrt{\langle x,x\rangle_M/D}.
\]
For perturbed data suppose the full equation
$z+e_z=(F_0+[e_x\ e_y])c+r$ holds, where
$\|e_x\|_M\le\epsilon_x$, $\|e_y\|_M\le\epsilon_y$,
$\|e_z\|_M\le\epsilon_z$, $\|r\|_M\le\rho$ and all budgets are
nonnegative. If $q=\epsilon_x\ell_\alpha+\epsilon_y\ell_\beta<1$, put
\[
 B=\frac{\epsilon_z+\rho+\epsilon_x|\alpha_0|
                    +\epsilon_y|\beta_0|}{1-q}.
\]
Then $|\alpha-\alpha_0|\le\ell_\alpha B$ and
$|\beta-\beta_0|\le\ell_\beta B$, and the perturbed columns remain
independent. If these bounds give $0<a_L\le\alpha\le a_U$ and
$0<b_L\le-\beta\le b_U$, define the associated positive durations by
\eqref{eq:source-scalar-clock-duration}. They obey
\begin{align}
 \frac{2b_L}{a_U(a_U+1)}&\le h_{j-1}^2
       \le\frac{2b_U}{a_L(a_L+1)},\notag\\
 \frac{2b_La_L}{a_L+1}&\le h_j^2
       \le\frac{2b_Ua_U}{a_U+1}.
 \label{eq:source-scalar-clock-stable-duration}
\end{align}
\end{proposition}
\begin{proof}
The two rows of $(F_0^*F_0)^{-1}F_0^*$ have the stated norm bounds.
Applying them to the perturbed equation bounds each coefficient error by
its $\ell$ times $b=\epsilon_z+\rho+\epsilon_x|\alpha|+
\epsilon_y|\beta|$. Substitution gives $b\le(1-q)B+qb$, hence $b\le B$.
In the coefficient norm $\max(|c_1|/\ell_\alpha,|c_2|/\ell_\beta)$,
the left-inverse times $[e_x\ e_y]$ has norm at most $q<1$; a vector
in the perturbed kernel must vanish. Formula
\eqref{eq:source-scalar-clock-stable-duration} follows by monotonicity
from \eqref{eq:source-scalar-clock-duration}.
\end{proof}

For the actual record, $\|\mathsf A\|_M<614$. A hypothetical error at most
$\epsilon=10^{-8}$ in each complete configuration therefore permits
$\epsilon_x=\epsilon_z=2\epsilon$, $\epsilon_y=614\epsilon$ and
$\rho=0$ for a hidden exactly configuration-stationary history. Every
$q<7.64\,10^{-6}$ and the coefficient sign margins stay positive.
Intersecting bounds on shared intervals and summing all $21$ durations gives,
in either schedule, the exact rational enclosure
\[
 1.341635903357\le t_{21}-t_0\le1.341645669638.
\]
These are conditional stability bounds, not measured noise. Nonempty
interval intersections do not by themselves establish existence of an
arbitrary noisy stationary history. A nonzero residual of the configuration
Euler equation must be multiplied by $h_j$ before it is bounded by $\rho$;
the residual units cannot be identified without that factor.

The code independently authenticates all $5{,}888$ events and $34{,}944$
reads, reconstructs the graph layers, and checks the Gram and stability
calculations with exact rational and quadratic-field arithmetic. The same
reconstructed values agree across the two schedules. This finite action
clock selects neither physical seconds nor native dynamics, regional
time-slice generation, a count-volume interpretation or a quantum clock.
Resource pairing, the action, its quadrature, orientation and relative
normalization remain explicit inputs.
